\documentclass[10pt]{article}
\usepackage{icmlko}

\icmlmeta{IP 파운데이션 월드모델}{84}{덮개율 0.6이 가장 센 축을 버리고 있었다}{새 물리량은 있었다 --- 못 받은 것은 구조였다}

\begin{document}

\twocolumn[
\icmlheader

\begin{abstract}
노트 83이 ``새 물리량이 유일하게 안 열어 본 문''이라 했다. 셋을 만들어 열었다.
\textbf{경쟁 밀도}(자기보다 먼저 같은 창에 나온 것의 수)는 공통 축으로도
고유 축으로도 $-$0.02다. \textbf{전형성}(태그가 도메인 전형에 얼마나 가까운가)은
$-$0.0085다. 둘 다 실패. \textbf{직전작 성적}(같은 창작자의 자기보다 먼저
나온 작품의 라벨)은 다르다 --- 도메인 안 상관이 모바일 \textbf{$+$0.709},
펀딩 $+$0.649, 게임 $+$0.629다. 모바일의 $+$0.709는 그 도메인 인자 공간의
자기 상관($+$0.656)\emph{보다 높다}. \textbf{축 하나가 여섯 축을 이긴다.}
그런데 모형에 넣으면 $\Delta-$0.0008이고, \textbf{여덟 도메인 중 일곱에서
정확히 $\pm$0.0000이다.} 희석이 아니라 \textbf{완전 무시}다. 범인을 찾았다 ---
\texttt{factor\_space} 의 \texttt{min\_cov=0.6}. 직전작 성적의 덮개율이
12$\sim$58\%라 전부 문턱 아래로 버려지고, 92\%인 세계애니만 통과해 그것만
$+$0.0279였다. \textbf{관측된 레코드에서만 재면} 모바일 자기 $\rho$가 $+$0.665에서
\textbf{$+$0.784}로, 애니가 $+$0.324에서 $+$0.426으로 오른다(일곱 중 다섯 상승).
문턱이 있는 이유도 분명하다 --- 인자 공간이 \textbf{완전 케이스만} 쓰므로
덮개율 36\%인 축을 넣으면 표본이 36\%로 준다. \textbf{물리량은 있었고 구조가
못 받았다.} 노트 83의 로드맵을 고친다 --- 남은 지렛대는 물리량이 아니라
\textbf{구조}이고, 그것이 노트 60--62의 신경망 질문을 다시 연다.
\end{abstract}
]

\begin{figure*}[t]
\centering
\includegraphics[width=\textwidth]{figs/gate.pdf}
\caption{왼쪽 --- 직전작 성적이 관측된 레코드에서만 잰 자기 $\rho$.
오른쪽 --- 새 물리량 셋의 판정치 효과.}
\end{figure*}

\section{새 물리량 셋}

\begin{table}[t]
\centering\tsize
\caption{판정치 효과. 현행 $+$0.4038.}
\begin{tabular}{lr}
\toprule
경쟁 밀도(공통 축, 창 365일) & $-$0.0195 \\
경쟁 밀도(고유 축) & $-$0.0218 \\
전형성(고유 축) & $-$0.0085 \\
직전작 성적(고유 축) & $-$0.0008 \\
\bottomrule
\end{tabular}
\end{table}

\textbf{경쟁 밀도}는 열 도메인이 문자 그대로 같은 것을 재는 첫 양이라 네 번째
공통 축 후보였다(노트 37: 공통 축이 성분보다 적으면 회전이 과소결정된다).
창을 90일에서 1,825일까지 바꿔도 전부 $-$0.017$\sim-$0.024다. 축이 이미 시간으로
탈추세돼 있어 시간과 거의 같은 양을 다시 넣는 셈이기 때문으로 보인다.

\section{직전작 성적은 다르다}

\begin{table}[t]
\centering\tsize
\caption{직전작 성적과 라벨의 도메인 안 순위 상관.}
\begin{tabular}{lrr}
\toprule
& $\rho$ & 관측 \\
\midrule
\textbf{모바일} & \textbf{$+$0.709} & 722 \\
펀딩 & $+$0.649 & 48 \\
게임 & $+$0.629 & 184 \\
애니 & $+$0.365 & 1204 \\
도서 & $+$0.337 & 120 \\
웹툰 & $+$0.305 & 1104 \\
세계애니 & $+$0.085 & 2723 \\
만화 & $+$0.049 & 624 \\
\bottomrule
\end{tabular}
\end{table}

지금 매장 노출도는 창작자의 사전 작품 \emph{수}만 센다. 그 작품들이 \emph{얼마나
잘됐는지}는 안 쓴다. 같은 창작자의 먼저 나온 작품의 라벨은 출시 시점에 이미
알려진 값이라 누출이 아니고, 제품 관점에서도 맞다 --- 우리는 우리 과거 팝업의
성적을 안다.

\claimx{직전작 성적은 이 프로젝트에서 나온 가장 강한 단일 변수다. 모바일에서
$\rho=+$0.709로 그 도메인 인자 공간 전체의 자기 상관($+$0.656)보다 높다.}
{여덟 중 둘(만화 $+$0.049, 세계애니 $+$0.085)은 거의 0이다. 두 도메인 다
제작사/스튜디오를 키로 썼는데 대형 스튜디오가 작품 수가 많아 직전작이 대표성이
없기 때문일 수 있다.}

\section{그런데 모형이 안 받는다}

\begin{table}[t]
\centering\tsize
\caption{도메인 안 자기 $\rho$ --- 직전작 축을 더할 때.}
\begin{tabular}{lrrr}
\toprule
& 현행 & $+$직전작 & $\Delta$ \\
\midrule
게임 & $+$0.3645 & $+$0.3645 & \textbf{$\pm$0.0000} \\
도서 & $+$0.2016 & $+$0.2016 & \textbf{$\pm$0.0000} \\
펀딩 & $+$0.1899 & $+$0.1899 & \textbf{$\pm$0.0000} \\
웹툰 & $+$0.4803 & $+$0.4803 & \textbf{$\pm$0.0000} \\
애니 & $+$0.2913 & $+$0.2913 & \textbf{$\pm$0.0000} \\
모바일 & $+$0.6559 & $+$0.6559 & \textbf{$\pm$0.0000} \\
만화 & $+$0.4773 & $+$0.4773 & \textbf{$\pm$0.0000} \\
\textbf{세계애니} & $+$0.4182 & $+$0.4461 & $+$0.0279 \\
\bottomrule
\end{tabular}
\end{table}

\textbf{정확히 0이다.} 희석이면 작은 값이라도 나와야 한다. 세계애니 하나만
다르고, 그 하나가 유일하게 덮개율 92\%다.

\claimx{\texttt{factor\_space} 의 \texttt{min\_cov=0.6} 문턱이 직전작 성적을
통째로 버리고 있었다. 덮개율이 12$\sim$58\%라 여덟 중 일곱이 문턱 아래이고,
문턱을 넘은 하나만 $+$0.0279를 냈다.}{문턱에는 이유가 있다 --- 인자 공간은
완전 케이스만 쓰므로 덮개율 36\%인 축을 넣으면 표본이 36\%로 준다. 문턱을
그냥 내리면 표본이 무너진다.}

\section{관측된 곳에서는 크게 오른다}

\begin{table}[t]
\centering\tsize
\caption{직전작이 관측된 레코드에서만 잰 자기 $\rho$.}
\begin{tabular}{lrrrr}
\toprule
& 덮개율 & 없이 & 넣고 & $\Delta$ \\
\midrule
\textbf{모바일} & 36\% & $+$0.665 & \textbf{$+$0.784} & \textbf{$+$0.119} \\
애니 & 58\% & $+$0.324 & $+$0.426 & $+$0.102 \\
게임 & 38\% & $+$0.444 & $+$0.505 & $+$0.061 \\
웹툰 & 39\% & $+$0.481 & $+$0.516 & $+$0.035 \\
세계애니 & 92\% & $+$0.421 & $+$0.446 & $+$0.025 \\
도서 & 49\% & $+$0.215 & $+$0.203 & $-$0.012 \\
만화 & 35\% & $+$0.504 & $+$0.477 & $-$0.027 \\
\bottomrule
\end{tabular}
\end{table}

일곱 중 다섯이 오르고 모바일이 \textbf{$+$0.784}에 닿는다 --- 노트 83이 낸
라벨 천장 추정 0.72$\sim$0.75\emph{보다 높다}. 둘 중 하나다. 천장 추정이 낮거나
(웹툰/애니 쌍에서 쟀다), 모바일 라벨의 신뢰도가 그 둘보다 높거나.

\section{로드맵을 고친다}

노트 83은 지렛대 점검표에서 ``새 물리량''만 미탐색으로 남기고 그것이 유일한
길이라 했다. \textbf{틀렸다.} 물리량은 있었고 --- 단일 상관 $+$0.709짜리가 ---
구조가 못 받았다.

\begin{quote}\fontsize{9.5}{11.5}\selectfont
\textbf{남은 지렛대는 구조다.} 인자 공간이 완전 케이스를 요구하는 한, 덮개율이
낮은 축은 아무리 세도 못 쓴다. 마스크를 \emph{입력으로} 받는 모형은 그 요구가
없다 --- 노트 60--62의 공유 인코더가 정확히 그 구조였고(값 다섯 $+$ 마스크
다섯), 그때는 PCA$+$프로크루스테스에 졌다. \textbf{그때는 못 쓸 축이 없었기
때문이다.} 지금은 있다.
\end{quote}

\section{한계}

\begin{itemize}[leftmargin=1.1em,itemsep=0.15em]
\item ``관측된 곳에서만'' 잰 자기 $\rho$는 표본이 다르므로 전체와 직접 비교
      못 한다. 부분집합이 유리한 쪽으로 치우쳤을 수 있다.
\item 직전작 성적은 \textbf{대상 라벨}(다른 레코드의)을 쓴다. 전이 설정에서
      대상에 쓰려면 그 도메인의 과거 라벨을 알아야 한다 --- 팝업은 알지만
      새 카테고리는 모른다.
\item 팝업에는 창작자/IP 키가 레코드에 없어 직전작 성적을 못 만들었다.
      제품에서 가장 값진 자리인데 못 쟀다.
\item 문턱을 실제로 내려서 재 보지 않았다. 표본이 얼마나 무너지는지 안 봤다.
\item 경쟁 밀도의 실패를 ``시간과 겹쳐서''로 설명했는데 확인 안 했다.
\end{itemize}

\section{다음}

\begin{enumerate}[leftmargin=1.3em,itemsep=0.15em]
\item \textbf{마스크를 입력으로 받는 모형을 다시 만든다}(노트 60--62의 재도전).
      이번에는 못 쓸 축이 있다는 점이 다르다.
\item 팝업 레코드에 IP/브랜드 키를 붙여 직전작 성적을 만든다. 제품에서 가장
      값진 자리다.
\item \texttt{min\_cov} 를 내리며 표본 손실과 이득의 교환을 잰다 --- 구조를
      안 바꾸고 얻을 수 있는 몫이 얼마인지.
\item 모바일 자기 $\rho$ $+$0.784가 천장 추정을 넘는 것 --- 천장을 다시 잰다.
\end{enumerate}

\vspace{0.6em}
\hrule height 0.3pt
\vspace{0.5em}
{\footnotesize
\textbf{참고문헌}\\[0.3em]
Little, R. J. A., Rubin, D. B. (2019). \emph{Statistical Analysis with Missing
Data}, 3rd ed. Wiley.\\[0.15em]
Josse, J. et al. (2019). On the consistency of supervised learning with missing
values. \emph{arXiv:1902.06931}.\\[0.15em]
Ilin, A., Raiko, T. (2010). Practical approaches to principal component analysis
in the presence of missing values. \emph{JMLR}, 11, 1957--2000.\\[0.15em]
Yoon, J. et al. (2018). GAIN: missing data imputation using generative
adversarial nets. \emph{ICML}, 5689--5698.\\[0.15em]
Ben-David, S. et al. (2010). A theory of learning from different domains.
\emph{Machine Learning}, 79(1--2), 151--175.
}

\end{document}
